%%%%%%%%%%%%%%%%%%%%%%%%%%%%%%%%%%%%%%%%%%%%%%%%%%%%%%%%%%%%%%%%%%%%%%
% Lớp 10 | Thống kê | Mô tả và biểu diễn dữ liệu trên các bảng và biểu đồ
% Dạng 0D6?2-3: Số liệu bất thường trên bảng số liệu
% Số câu: 7  (sinh tự động bởi tach_cau.py; từ nay thầy sửa trực tiếp file này)
%%%%%%%%%%%%%%%%%%%%%%%%%%%%%%%%%%%%%%%%%%%%%%%%%%%%%%%%%%%%%%%%%%%%%%

%%%%%%%%%%%%%%%% TRẮC NGHIỆM 4 PHƯƠNG ÁN %%%%%%%%%%%%%%%%

% nguồn: MaHoaNH_TracNghiem_Lop10
\begin{ex}%[0D6N2-3]%[ma:91398d76]
	Tâm ghi lại số liệu từ trang web của Tổng cục thống kê  bảng nhiệt độ không khí trung bình các tháng trong năm $2020$ tại một trạm quan trắc đặt ở thành phố Vinh.
	\begin{longtable}{|c|c|c|c|c|c|c|c|c|c|c|c|c|}
		\hline 
		Tháng&$1$&$2$&$3$&$4$&$5$&$6$&$7$&$8$&$9$&$10$&$11$&$12$\\ 
		\hline 
		Nhiệt  độ&$20{,}9$&$20{,}7$&$23{,}7$&$23$&$29{,}5$&$32{,}2$&$4{,}5$&$29{,}6$&$28{,}9$&$23{,}8$&$23{,}1$&$18{,}4$\\
		 \hline
	\end{longtable}\noindent
	Bạn Tâm đã ghi nhầm  nhiệt độ của một tháng trong bảng bên. Theo em bạn Tâm đã ghi nhầm số liệu của tháng mấy?
	\choice 
		 {Tháng $6$}
		{\True Tháng $7$} 
		{Tháng  $12$}
		{Tháng  $1$}
	\loigiai{ 
		 Tâm ghi nhầm nhiệt độ của tháng $7$.\\
		Do tháng $7$ là  vào mùa hè nên nhiệt độ trung bình trong tháng đó ở thành phố Vinh phải cao hơn $4{,}5^\circ$.
	} 
\end{ex}

% nguồn: Ok_TracNghiemLop11_Lan1
\begin{ex}%[0D6B2-3]%[ma:0d525781]
	Biết $A,B,C$ là các góc của tam giác $ABC$, mệnh đề nào sau đây đúng. 
	\choice 
		{$\sin  \left(A+C\right)=-\sin B$}
		{\True $\cos \left(A+C\right)=-\cos B$} 
		 {$\tan \left(A+C\right)=\tan B$}
		 {$\cot \left(A+C\right)=\cot B$}
	\loigiai{Vì $A,B,C$ là  ba góc của một tam giác suy ra $A+C=\pi-B.$\\
		Khi đó $\sin \left(A+C\right)=\sin \left(\pi-B\right)=\sin B;\cos \left(A+C\right)=\cos \left(\pi-B\right)=-\cos B.$\\ 
		$\tan \left(A+C\right)=\tan \left(\pi-B\right)=-\tan B;\cot \left(A+C\right)=\cot \left(\pi-B\right)=-\cot B.$ 
	} 
\end{ex}

% nguồn: Ok_TracNghiemLop11_Lan1
\begin{ex}%[0D6B2-3]%[ma:a6711304]
	Với mọi số thực $\alpha $, ta có $\sin \left(\dfrac{9\pi}{2}+\alpha\right)$ bằng 
	\choice 
		{$-\sin \alpha$} 
		{\True $\cos \alpha$} 
		{$\sin \alpha$} 
		 {$-\cos \alpha$}
	\loigiai{Ta có $\sin  \left(\dfrac{9\pi}{2}+\alpha\right)=\sin \left(4\pi+\dfrac{\pi}{2}+\alpha\right)=\sin \left(\dfrac{\pi}{2}+\alpha\right)=\cos \alpha.$ }
\end{ex}

% nguồn: Ok_TracNghiemLop11_Lan1
\begin{ex}%[0D6B2-3]%[ma:f23e7ce9]
	 Với mọi $\alpha \in \mathbb{R}$ thì $\tan \left(2017\pi+\alpha\right)$ bằng
	 \choice
		{$-\tan  \alpha$}
		{$\cot \alpha$} 
		 {\True $\tan \alpha$}
		 {$-\cot \alpha$}
	\loigiai{ Ta có $\tan  \left(2017\pi+\alpha\right)=\tan \alpha.$}
\end{ex}


%%%%%%%%%%%%%%%% TRẢ LỜI NGẮN %%%%%%%%%%%%%%%%

% nguồn: MaHoaNH_TraLoiNgan_Lop10
\begin{bt}%[0D6H2-3]%[ma:e4130295]
	 Một lớp gồm $40$ học sinh chia đều thành $4$ tổ. Nhân dịp Tết Nguyên Đán, các em cùng đóng góp để giúp đỡ các gia đình khó khăn cùng đón cái Tết đầm ấm. Biết rằng mỗi em đóng góp $10$ đến $15$ ngàn đồng. Cuối đợt đóng góp, lớp trưởng thống kê lại như bảng sau
	\begin{center}
		\begin{tabular}{|c|c|c|c|c|}
			\hline Tổ& $1$ & $2$ &  $3$ & $4$ \\
			\hline Số tiền (đồng) & $90\,000$ & $102\,000$ & $200\,000$ & $140\,000$ \\ 
			 \hline
		\end{tabular}
	\end{center}
	Hỏi  lớp trưởng thống kê sai mấy tổ?
	\shortans{2} 
	 \loigiai{
		Số học sinh mỗi tổ là $\dfrac{40}{4}=10$ học  sinh.\\
		Số tiền tối thiểu mỗi tổ quyên góp  là $10\cdot 10\,000=100\,000$\,đồng.\\
		Số tiền tối đa mỗi  tổ quyên góp là $10\cdot 15\ 000=150\,000$\,đồng.\\
	 Vậy lớp trưởng thống kê sai $2$ tổ (tổ $1$ và tổ $3$).}
\end{bt}


%%%%%%%%%%%%%%%% TỰ LUẬN %%%%%%%%%%%%%%%%

% nguồn: Ok_TuanLuan_Toan11_Lan1
\begin{bt}%[0D6K2-3]%[ma:042c9c35]
	Tính giá trị của biểu thức $B=\cos 20^\circ+\cos 40^\circ+\cos 60^\circ  +\cdots+\cos180^\circ$.
	\loigiai{ 
		Ta có  $\cos\left(180^\circ-x\right)=-\cos x\Rightarrow \cos x+\cos\left(180^\circ-x\right)=0$.\\
		$\Rightarrow \heva{&\cos20^\circ+\cos160^\circ=0\\& \cos 40^\circ+\cos140^\circ=0\\&\cos  60^\circ+\cos 120^\circ=0\\& \cos 80^\circ+\cos 100^\circ=0}\Rightarrow B=4.0 +\cos 180^\circ =0-1=-1$.
	 }
\end{bt}

% nguồn: Ok_TuanLuan_Toan11_Lan2
\begin{ex}%[0D6K2-3]%[ma:e6b34706]
	Tính giá trị của biểu thức $B=\cos  20^\circ+\cos 40^\circ+\cos 60^\circ +\cdots+\cos180^\circ$.
	 \loigiai{
		Ta có $\cos\left(180^\circ-x\right)=-\cos x\Rightarrow \cos  x+\cos\left(180^\circ-x\right)=0$.\\
		$\Rightarrow \heva{&\cos20^\circ+\cos160^\circ=0\\& \cos 40^\circ+\cos140^\circ=0\\&\cos 60^\circ+\cos 120^\circ=0\\& \cos 80^\circ+\cos 100^\circ=0}\Rightarrow B=4.0 +\cos  180^\circ =0-1=-1$.
	} 
\end{ex}
